% !PRINTED 742
\weberpage{742}

\S.\,202.

\gesperrt{Positive Units.}

The results last obtained show that the system of independent
units \S.\,201, (7)
\[ \varepsilon'_1, \varepsilon'_2, \ldots, \varepsilon'_{\nu'-1}, \; \mathcal{E}_0, \mathcal{E}_1, \ldots, \mathcal{E}_{\nu'-1} \]
forms no \gesperrt{fundamental system} in the field $H_m$. For otherwise the
% [sic] -- the source omits a separator between the coefficient families $m_{k,i}$ and $M_{k,i}$; the printed form is preserved.
numbers $m_{k,i} M_{k,i}$ in the formulae \S.\,201, (13) would all have
to be divisible by 2, and the determinant $M$, whose value we
found $= 2^{\nu'}$, would be divisible by $2^{\nu-1}$.

In place of the system of equations (18), \S.\,201, we now consider
the following congruences:
\weq{1}{
\begin{aligned}
\sum_{1,\,\nu-1}^{i} b_i m_{1,i} &\equiv 1 \\
\sum_{1,\,\nu-1}^{i} b_i m_{2,i} &\equiv 0 \pmod{2}, \\
& \cdots\cdots\cdots\cdots \\
\sum_{1,\,\nu-1}^{i} b_i m_{\nu'-1,i} &\equiv 0
\end{aligned}
}
and we prove that it is possible to satisfy them by integral
values of the $b_i$.

If it were not possible, then from the $\nu'-2$ congruences
\[ \sum_{1,\,\nu-1}^{i} b_i m_{k,i} \equiv 0 \pmod{2} \quad k = 2, 3, \ldots, \nu'-1 \]
the last one
\[ \sum_{1,\,\nu-1}^{i} b_i m_{1,i} \equiv 0 \pmod{2} \]
would have to follow as a \gesperrt{consequence}, and then, as in the
preceding section, it would follow further:
\[ \sum_{1,\,\nu-1}^{i} b_i M_{k,i} \equiv 0 \pmod{2} \quad k = 0, 1, \ldots, \nu'-1. \]
